% 第一章 绪论

\section{绪论}

\subsection{研究背景与意义}
\par 智能家居作为物联网（IoT）在居民生活场景中的典型落地形态，正从“单品智能、各自为政”逐步走向“全屋互联、场景联动、自然交互”。一方面，传感器、嵌入式与无线通信的成本持续下降，使得对温湿度、气压、空气质量等环境要素的低功耗连续监测成为可能；另一方面，\textbf{标准化与互操作}正在重塑产业格局。以Matter为代表的连接标准在近年持续演进，引入多生态系统协同接入、家庭网络基础设施（如Thread边界路由能力）与能源管理扩展等能力，推动智能家居从“厂商生态壁垒”向“IP化互联底座”迁移\cite{csa_matter14}。随着互联“管道”逐渐打通，行业竞争焦点进一步上移到\textbf{边缘智能与家庭AI}层面，即在本地完成更低时延、更强隐私保护的感知理解与自动化决策\cite{eet_matter2026}。
\par 同时，从社会与时事角度看，居民对\textbf{居住安全、健康与能耗管理}的关注度不断提高：极端天气频发带来的室内舒适与安全问题、老龄化背景下的居家照护需求、以及“双碳”目标下的家庭能源精细化管理，都对智能家居系统提出了更高要求。智能家居不再只是“远程开关控制”，而是需要围绕“多源感知—可靠连接—数据服务—人机交互—安全与隐私”形成闭环能力。
\par 当前智能家居行业发展现状呈现“多协议并存、标准加速统一、边缘智能上移”的趋势。一方面Wi-Fi/蓝牙/Zigbee/Thread等链路长期共存；另一方面，Matter等标准通过统一数据模型与互操作机制，降低跨品牌部署门槛，并将“能源管理、设备状态上报、低功耗间歇连接”等能力纳入标准演进路线\cite{csa_matter14}。
\par 在智能家居场景中，物联网通常由端侧感知节点、家庭网关、云/本地服务与应用层组成。端侧完成多传感器采集与部分边缘处理，网关实现协议适配与数据汇聚，服务端提供数据存储、可视化与控制接口。MQTT凭借发布/订阅、轻量与弱网适应等优势，成为智能家居设备控制与状态同步的常用协议之一\cite{mqtt_oasis}。
\par 面向电子信息工程专业的综合训练需求，智能家居系统天然覆盖“传感器与信号采集、嵌入式软件设计、串口/网络通信协议、边缘视觉、人机交互、安全认证与工程实现”等关键能力。本文以“STM32采集 + ESP32网关 + OpenMV视觉 + Node.js平台 + LLM语音助手”的端边云协同方案为载体，完成从硬件到软件、从数据到交互的系统级设计与实现，对提升系统工程能力与落地实践具有现实意义。

\subsection{国内外研究现状}
\par 国内外智能家居研究与产品演进大体可归纳为四条主线：系统架构从云中心向端边云协同演进、环境监测由单点测量走向多源融合与长期稳定性、身份识别从外部安防延伸到室内访问控制、语音交互从固定指令走向大模型驱动的自然语言与工具调用式执行。
\par 在智能家居系统架构研究方面，传统方案多依赖云端平台与各自封闭的生态入口，设备互联与场景联动往往受限于厂商协议与网关能力。近年来，基于IP的互操作标准（如Matter）推动跨生态接入与统一数据模型；同时，为降低时延与提升隐私性，端侧与家庭网关承担更多本地处理任务，形成“端侧感知 + 网关汇聚 + 本地/云端服务”的分层架构\cite{csa_matter14,eet_matter2026}。
\par 在环境监测技术方面，家庭环境监测关注温湿度、气压、空气质量（VOC/可燃气体等）等指标。研究重点从“传感器选型与接口驱动”扩展到“标定与漂移补偿、采样策略与功耗控制、异常检测与数据质量保障”。在工程实现上，I\(^2\)C等总线传感器便于集成，但在长期运行、供电噪声与通信错误等条件下仍需完善的容错机制与数据校验策略。
\par 在人脸识别应用方面，人脸识别在门禁与室内权限控制中具有自然、无接触的优势，但也带来隐私与数据安全挑战。为减少敏感图像上云，端侧视觉与边缘推理成为重要方向。OpenMV等微型视觉平台能够在资源受限条件下完成基于Haar级联的人脸检测与基于LBP等特征的轻量识别，适合在教学与原型系统中快速验证“端侧识别 + 控制联动”的可行性。
\par 在语音交互技术方面，早期语音控制主要依赖固定唤醒词与有限指令集，扩展性受限。近两年，大语言模型（LLM）在“意图理解、对话管理、工具调用”方面能力提升，使语音助手逐步从“语音转文本 + 规则解析”转向“ASR—LLM—TTS”流水线，并通过结构化工具调用将自然语言映射为受控的系统操作，从而在可用性与可控性之间取得更好的平衡。

\subsection{研究内容与目标}
\par 本文面向家庭环境监测与便捷交互需求，设计并实现一套可运行的智能家居原型系统，覆盖端侧采集、网关传输、服务器存储与可视化、人脸识别与语音交互等功能模块。系统总体思路为：STM32侧负责多传感器采集与结构化上报；ESP32-S3侧作为家庭网关完成数据融合、显示、联网与消息转发；OpenMV侧完成端侧人脸采集与识别并提供图像回传能力；服务器侧基于Node.js实现数据接收、用户认证、历史数据管理、可视化与基于MQTT的命令下发；语音助手侧通过LLM工具调用与服务器接口联动，实现自然语言查询与控制。
\par 系统功能需求主要包括环境参数采集与上报（温度、湿度、气压、VOC/气体相关指标等），并在网关端实现本地信息展示（TFT屏显示关键指标与网络状态）。同时，系统需要支持人脸系统的采集、识别、状态查询、名单管理与图像获取等操作，并可通过消息机制与服务器/网页端联动；Web端需提供登录鉴权、历史数据查询、图表展示与系统设置等功能；语音交互部分需支持自然语言查询与控制，并通过工具调用触发受控接口完成任务执行。
\par 技术指标主要体现在采集与刷新、端到端时延、稳定性与安全性基础等方面。在采集与刷新方面，传感器采样周期应可配置，显示与上传采用相互独立的定时策略以避免阻塞；在端到端时延方面，从端侧采集到服务器入库/可视化展示应具备秒级更新能力，控制指令链路（网页/服务器—MQTT—网关—OpenMV）应具备可感知的实时响应；在稳定性方面，应支持网络断开重连、异常数据处理与长时间运行下的数据持续写入与服务可用性；在安全性基础方面，服务器端需具备用户注册/登录与会话校验机制，控制接口默认在认证后访问。
\par 创新点主要体现在三方面。本文采用STM32+ESP32分层协作，将多传感器采集与网络/显示/协议栈解耦，通过串口进行结构化数据上送，以降低单节点负载并提升系统可维护性；在视觉侧，基于OpenMV实现端侧人脸采集/识别与图像获取的串口命令协议，并通过MQTT与服务器API实现跨端协同控制；在人机交互侧，采用“ASR—LLM—TTS”流水线，通过工具定义与参数校验把自然语言意图映射为受控接口调用，提升交互自然性与执行可控性。

\par 配图建议如下。图1-1建议为“系统整体架构示意图”，画法为三层架构：顶层为应用层（Web前端、语音助手），中间层为通信层（HTTP、MQTT、UART），底层为硬件层（STM32采集、ESP32网关、OpenMV视觉、TFT显示），并用箭头标注数据流与控制流方向。该图可用Visio/diagrams.net绘制，也可用PPT绘制后导出。
\begin{figure}[H]
\centering
\includegraphics[width=0.92\textwidth]{paper/fig/1-1.png}
\caption{系统整体架构示意图}
\label{fig:system_overview}
\end{figure}

% \subsection{论文组织结构}
% \par 本文内容组织如下。
% \par 第1章绪论介绍研究背景与意义、国内外研究现状、本文研究内容与目标以及论文结构安排。第2章给出系统需求分析、总体架构与模块划分，并说明硬件平台、通信协议与软件框架的技术选型依据。第3章阐述STM32采集端、ESP32网关端与OpenMV视觉端的硬件组成、接口连接与关键实现思路。第4章描述UART数据帧设计、网络通信与MQTT主题/消息格式，以及异常处理与重连机制。第5章介绍Node.js服务器架构、数据与用户管理、Web前端可视化以及设备接入与控制接口。第6章说明ASR、LLM、TTS与工具模块的组成、提示词与工具调用机制及与服务器接口的集成方法。第7章给出软硬件实现过程、联调方法与测试方案，并对功能、性能与稳定性测试结果进行分析。第8章总结本文工作与创新点，分析不足并展望后续优化方向。
